Let $\xi$ be a non-zero real number and let $n$ be a positive integer. Dirichlet's theorem (1842) is one of the most basic results of Diophantine approximation. It shows that for any real number $H > 1$, there exists a non-zero integer point $(x_0,\dots,x_n)\in\bZ^{n+1}$ such that
\begin{equation}
\label{eq: Dirichlet system}
\max\big\{|x_1|,\dots,|x_n|\big\} \leq H \AND |x_0 + x_1\xi + \cdots + x_n\xi^n| \leq H^{-n}.
\end{equation}
It is natural to ask if we can improve the exponent $n$ of $H^{-n}$, and this question gives rise to two Diophantine exponents. The so-called \emph{uniform exponent of approximation} $\homega_n(\xi)$ (\resp the \emph{ordinary} exponent $\omega_n(\xi)$), is the supremum of the real numbers $\omega > 0$ such that the system
\begin{equation*}
\normH{P} \leq H \AND 0 <|P(\xi)| \leq H^{-\omega}
\end{equation*}
admits a non-zero solution $P\in\bZ[X]$ of degree at most $n$ for each sufficiently large $H$ (\resp for arbitrarily large $H$). Here, $\normH{P}$ denotes the (naive) \emph{height} of $P$, defined as the largest absolute value of its coefficients. These quantities have been extensively studied over the past half-century, see for example \cite{bugeaud2015exponents} for a nice overview of the subject. By Dirichlet's theorem, if $\xi$ is not an algebraic number of degree $\leq n$, then we have
\[
\omega_n(\xi) \geq \homega_n(\xi) \geq n,
\]
and it is well known that those inequalities are equalities for almost all real numbers~$\xi$ (with respect to the Lebesgue measure). Note that if $\xi$ is an algebraic number of degree $d$, then $\homega_n(\xi)$ and $\omega_n(\xi)$ are both equal to $\min\{n,d-1\}$ (it is a consequence of Schmidt's subspace theorem, see \cite[Th.\,2.10]{bugeaud2015exponents}). We can therefore restrict our study to the set of transcendental real numbers. The initial question ``can we improve the exponent $n$ in Dirichlet's theorem?'' may be rephrased as follows: ``does there exist a transcendental real number $\xi$ satisfying $\homega_n(\xi) > n$?''. For $n=1$ the answer is negative and rather elementary to prove, so the first non-trivial case is $n=2$. Before the early $2000$s, it was conjectured that no such number existed. This belief was swept away by Roy's extremal numbers \cite{roy2003approximation}, \cite{roy2004approximation}, \cite{arbourRoyCriterionDegreTwo}, whose exponent $\homega_2$ is equal to the maximal possible value $(3+\sqrt 5)/2 = 2.618\cdots$. Since then, several families of transcendental real numbers whose uniform exponent $\homega_2$ is greater than $2$ have been discovered (see for example \cite{roy2007two}, \cite{bugeaud2005exponentsSturmian}, \cite{poels2017exponents, poelsExpoGeneralClass2021}). However, for $n \geq 3$ the mystery remains, and it is still an open question whether or not there exists $\xi\in\bR\setminus\overline{\bQ}$ with $\homega_n(\xi) > n$.
